% Generated draft: requires mathematical and source review.
\documentclass[11pt]{article}
\usepackage[margin=25mm]{geometry}
\usepackage{fontspec}
\setmainfont{Latin Modern Roman}
\usepackage{amsmath,amssymb,mathtools,unicode-math}
\setmathfont{Latin Modern Math}
\usepackage{enumitem,needspace,xurl,hyperref}
\setlength{\parindent}{0pt}
\setlength{\parskip}{5pt}
\setlength{\emergencystretch}{3em}
\newcommand{\N}{\mathbb{N}}
\newcommand{\Z}{\mathbb{Z}}
\newcommand{\Q}{\mathbb{Q}}
\newcommand{\R}{\mathbb{R}}
\newcommand{\C}{\mathbb{C}}
\newcommand{\F}{\mathbb{F}}
\newcommand{\A}{\mathbb{A}}
\newcommand{\PP}{\mathbb{P}}
\newcommand{\E}{\mathbb{E}}
\newcommand{\eps}{\varepsilon}
\newcommand{\dd}{\,\mathrm{d}}
\begin{document}
\section*{UnsolvedMath ID: 30005408}
\textbf{Source problem code:} catalogue entry
\subsection{Definitions and mathematical statement}
Let \(k\) be the coefficient field and let \(R=k[[x,y]]\) be the ring of formal power series in two variables. The supplied record imposes no characteristic or algebraic-closure hypothesis on \(k\). A monomial ideal is an ideal generated by monomials \(x^a y^b\), with \(a,b\in\Z_{\geq 0}\). Here a zero-dimensional monomial ideal \(E\subset R\) is understood to have finite colength \(n=\dim_k(R/E)\).

The record associates to \(E\) a canonical Hilbert--Burch matrix \(H\), a space \(T'(E)\) of admissible perturbation matrices, and the Gr\"obner cell \(V(E)\) of ideals with initial ideal \(E\). For a matrix \(M\), write \(I_t(M)\) for the ideal generated by its \(t\times t\) minors. In the intended Hilbert--Burch presentation, these are the maximal minors of \(H+N\). The record does not specify the size of \(H\), the value of \(t\), the construction and normalization of \(H\), the admissibility conditions defining \(T'(E)\), or the order and initial-ideal convention defining \(V(E)\). Thus these objects retain their source-dependent meanings; the supplied information does not give a self-contained definition of the proposed map.

\textbf{Conjectural target, subject to those unresolved definitions.} For every zero-dimensional monomial ideal \(E\subset k[[x,y]]\) in the intended scope, the assignment
\[
\Phi_E:T'(E)\longrightarrow V(E),\qquad N\longmapsto I_t(H+N),
\]
is a well-defined bijection. Equivalently, every admissible \(N\) produces an ideal in \(V(E)\), and
\[
\forall J\in V(E)\quad \exists!\,N\in T'(E)\quad J=I_t(H+N).
\]
The scope over coefficient fields, and whether the unit ideal is excluded, require clarification from the source.

The supplied original excerpt reports that the conjecture holds for lex-segment ideals, an unspecified slightly more general class, and complete intersections. It also reports agreement, for \(n\leq 30\), of the cell dimensions predicted by the conjecture for a cellular decomposition of \(\operatorname{Hilb}^{n}(k[[x,y]])\). This dimension agreement is reported evidence, not a proof of bijectivity. The supplied dated assessment labels the general problem open and reports injectivity and surjectivity as unproved; these are supplied source assessments, not independently verified conclusions.
\subsection{Short English statement}
The proposed parametrization says that every ideal in the Gr\"obner cell of a finite-colength monomial ideal arises from exactly one admissible perturbation of its canonical Hilbert--Burch matrix. Precise matrix conventions, admissibility conditions, initial-ideal conventions, and field hypotheses are missing from the supplied record, so this remains a source-dependent, statement-only draft requiring mathematical review.
\subsection{Sources}
\begin{itemize}
\item Local UnsolvedMath registry, supplied complete source record: permanent ID \(30005408\), title ``Gr\"obner-Cell Parametrization of Punctual Hilbert Schemes,'' exact source problem code \texttt{catalogue entry}. The record gives proposed year 2023 and names no proposer.
\item \emph{Resolutions in Local Algebra and Singularity Theory} (2023), Oberwolfach Reports, pp.~327--395, DOI \texttt{10.4171/owr/2023/6}: \url{https://doi.org/10.4171/owr/2023/6}. The original excerpt includes the report page marker 347 and supplies the special-case and dimension evidence described above.
\item Supplied literature reference: ``OWR-2023-6, Conjecture 6,'' \url{https://oa.tib.eu/renate/server/api/core/bitstreams/11773ea5-b1e5-4e9e-b650-808e9d9dc944/content}.
\item Supplied literature reference: ``Homs--Winz, canonical Hilbert--Burch matrices,'' \url{https://doi.org/10.1016/j.jalgebra.2021.04.021}.
\item Literature assessment dated exactly \texttt{2026-08-22}: the supplied curation labels the problem open, reports injectivity and surjectivity as unproved in general, and treats the checks for \(n\leq 30\) as evidence only. This is supplied dated curation, not independently verified literature assessment. The record labels its statement \texttt{corrected\_verified} and reports direct checking of the official report; that verification claim is also supplied metadata and has not been independently checked here.
\end{itemize}
\end{document}
